\documentclass[
 pdftex, 
 oneside, % Einseitiger Druck.
 12pt, % Schriftgroesse
 parskip=half, % Halbe Zeile Abstand zwischen Absätzen.
 headsepline, % Linie nach Kopfzeile.
 %footsepline, % Linie vor Fusszeile.
 abstracton % Abstract Überschriften
%  ngerman % Translator
]{scrreprt}  
 
%Seitenränder:
\usepackage[ 
	top=25mm, 
	left=30mm, 
	right=25mm, 
	bottom=32mm,
	headsep=10mm, 
	footskip=15mm
	]{geometry} 
	
\usepackage{lscape}
	

%Seite von oben bis unten füllen und ggf. Abstände minimalst dehnen	

%\usepackage{fullpage} %Seitengroesse
\usepackage[activate]{microtype} %Zeilenumbruch und mehr

% Zeichencodierung:
\usepackage[T1]{fontenc}
\usepackage[utf8]{inputenc}
  
\usepackage[onehalfspacing]{setspace} % Zeilenabstand
\usepackage{makeidx} % Index-Erstellung
\usepackage[english]{babel} % Lokalisierung (neue deutsche Rechtschreibung)
\usepackage[babel,german=quotes]{csquotes} % deutsche Anführungszeichen

\usepackage{graphicx} % Grafiken
%Grafiken immer im selben Kapitel einfügen wo sie referenziert werden
\usepackage{placeins} 

% einfachere Querverweise mit Seitenangabe usw.:
\usepackage{prettyref}
\newrefformat{cha}{Kapitel~\ref{#1} auf Seite~\pageref{#1}}  % Für Kapitel 
\newrefformat{sec}{Abschnitt~\ref{#1} auf Seite~\pageref{#1}} % Für Abschnitte
\newrefformat{fig}{Abbildung~\ref{#1} auf Seite~\pageref{#1}}% Für Abbildungen
\newrefformat{tab}{Tabelle~\ref{#1} auf Seite~\pageref{#1}}	 % Für Tabellen
\newrefformat{lst}{Quellcode~\ref{#1} auf Seite~\pageref{#1}} % für Quellcodeliings 
\newrefformat{uc}{auf Seite~\pageref{#1}} % für Use-Cases 
 
%merhspaltige Elemente möglich machen
\usepackage{multicol} 

%Mathe-Zeugs und Text-Anordnung
\usepackage{amsmath}
\usepackage{amssymb} 

%serifenlose Schriftart:
\renewcommand{\familydefault}{\sfdefault}
\usepackage{sfmath}

\usepackage{float}
\restylefloat{table}
   
% Spezielle Tabellenform fuer Deckblatt
\usepackage{longtable}
\setlength{\tabcolsep}{10pt} %Abstand zwischen Spalten
\renewcommand{\arraystretch}{1.5} %Zeilenabstand

% Farben
\usepackage{color}
\definecolor{LinkColor}{rgb}{0,0,0.5}
\definecolor{ListingBackground}{rgb}{0.92,0.92,0.92}

%Farben in Tabellen nutzen
\usepackage{colortbl}
\definecolor{light-gray}{gray}{0.80}

\clubpenalty=10000
\widowpenalty=10000
\displaywidowpenalty=10000 

% Quellcode
\usepackage{listings}
\usepackage{scrhack} %bessere Integration in KOMA-Script, Warning wegbekommen
\definecolor{javared}{rgb}{0.6,0,0} % for strings
\definecolor{javagreen}{rgb}{0.25,0.5,0.35} % comments
\definecolor{javapurple}{rgb}{0.5,0,0.35} % keywords
\definecolor{javadocblue}{rgb}{0.25,0.35,0.75} % javadoc

\lstset{
	language=Java, % Sprache des Quellcodes
	numbers=left, % Zelennummern links
	stepnumber=1, % Jede Zeile nummerieren.
	numbersep=5pt, % 5pt Abstand zum Quellcode
	lineskip=1pt, % Zeilenabstand im Quellcode
	numberstyle=\tiny, % Zeichengrösse 'tiny' für die Nummern.
	breaklines=true, % Zeilen umbrechen wenn notwendig.
	breakautoindent=true, % Nach dem Zeilenumbruch Zeile einrücken.
	postbreak=\space, % Bei Leerzeichen umbrechen.
	tabsize=2, % Tabulatorgrösse 2
	basicstyle=\ttfamily\footnotesize, % Nichtproportionale Schrift, klein für den Quellcode
	showspaces=false, % Leerzeichen nicht anzeigen.
	showstringspaces=false, % Leerzeichen auch in Strings ('') nicht anzeigen.
	%basicstyle=\ttfamily, 
	morekeywords={var, function},
	keywordstyle=\color{javapurple}\bfseries,
	stringstyle=\color{javared},
	commentstyle=\color{javagreen},
	morecomment=[s][\color{javadocblue}]{/**}{*/},
	backgroundcolor=\color{ListingBackground}, % Hintergrundfarbe des Quellcodes
	% setzen.
	captionpos=b, % sets the caption-position to bottom
	%extendedchars=true, % Alle Zeichen vom Latin1 Zeichensatz anzeigen.
    literate={ö}{{\"o}}1 %deutsche Sonderzeichen auch in Listings
    		 {Ö}{{\"O}}1 %funktionierend machen durch escapen
         	 {Ü}{{\"U}}1
         	 {ü}{{\"u}}1
         	 {Ä}{{\"A}}1
         	 {ä}{{\"a}}1
         	 {ß}{{\ss}}1
}
\renewcommand{\lstlistlistingname}{Quellcodeverzeichnis}
\renewcommand{\lstlistingname}{Quellcode}

%Glossar-Einträge konfigurieren:
\usepackage[
	nonumberlist, %keine Seitenzahlen anzeigen
	%acronym,      %ein separates Abkürzungsverzeichnis erstellen
	toc,          %Einträge im Inhaltsverzeichnis
	section]      %im Inhaltsverzeichnis auf section-Ebene erscheinen
{glossaries}
\glsenablehyper   %Abkürzungen ins Glossar verlinken
\renewcommand*{\glspostdescription}{} %Punkt nach Beschreibung deaktivieren

%Wort aus Glossar bei erster Verwendung hervorheben
\defglsdisplayfirst{\emph{#1}} 
\defglsdisplayfirst[acronym]{\emph{#1}} 
 
 
%Literaturverzeichnis formatieren
\bibliographystyle{plaindin} %Literaturverzeichnis nach DIN-Norm darstellen
  
% Verschiedene Schriftarten
%\usepackage{goudysans}
%\usepackage{lmodern}
%\usepackage{libertine}
\usepackage{palatino} 


\newcommand{\pdftitel}{\autor \space \arbeit \space \semester}
\newcommand{\pdfsubject}{\arbeit \space \semester}
\newcommand{\autor}{Ephraim Petry \& Nick Herrmannsdörfer}
\newcommand{\arbeit}{Studienarbeit}
\newcommand{\semester}{Semester 5 + 6}


\newcommand{\titel}{Entwicklung einer Lern-Feedback-Plattform für Vorlesungen}
\newcommand{\subtitel}{Lecture Monitoring - LeMon}
\newcommand{\martrikelnrEphraim}{2767400}
\newcommand{\martrikelnrNick}{1655361}
\newcommand{\kurs}{TINF11D}
\newcommand{\datumAbgabe}{06. Juni 2014}
\newcommand{\abschluss}{Bachelor of Science}
\newcommand{\studiengang}{Informatik}
\newcommand{\vertiefung}{Angewandte Informatik}
\newcommand{\dhbw}{Stuttgart} 
\newcommand{\betreuer}{Frau Dr. Barbara Dörsam}
\newcommand{\zeitraum}{zwei Semester}
\newcommand{\abgabeort}{Stuttgart}

% Titel, Autor und Datum
\title{\titel}
\author{\autor}
\date{\datum}
 
%URl-Umbrüche auch an Bindestrichen für lange URLs erlauben
 \usepackage[hyphens]{url}
% set colors for links and URLs und URLs nicht escapen
\usepackage[
	colorlinks=true, 
    urlcolor=blue,
    linkcolor=black,
    citecolor=black,
    filecolor=black,
    %menuecolor=black,
    pdftitle={\pdftitel},
	pdfauthor={\autor},
	pdfsubject={\pdfsubject},
	pdfcreator={pdflatex, LaTeX with KOMA-Script},
	pdfpagemode=UseOutlines, % Beim Oeffnen Inhaltsverzeichnis anzeigen
	pdfdisplaydoctitle=true, % Dokumenttitel statt Dateiname anzeigen.
	pdflang=de % Sprache des Dokuments.
	]{hyperref} % Hyperlinks als solche machen

% Fußnoten
\usepackage[
	perpage, %Fußnoten-Nummerierung auf jeder Swite neu beginnen
	hang,	 %Abstand zwischen Zahl und Text in der Fußnote 
	multiple, 
	stable]  %Fußn. in Überschriften anzeigen, aber nicht im Inhaltsverzeichnis
{footmisc}


\usepackage{fancyhdr}


 \pagestyle{fancy}{
 	 	\newcommand{\clearHeader}{\fancyhead{}\fancyfoot{}}
	\fancyhead{}
	\fancyfoot{}  	
 	\fancyfoot[C]{\thepage}
	\fancyhead[L]{\includegraphics[height=0.65cm]{images/dhbw-Logo.png}}
 	
 
 	\fancyhead[C]
 	{
 	 \begin{tabular}[b]{l}
 	  LeMon
 	 \end{tabular}
 	}

 %	\fancyhead[R]
 %	{
 %	 \begin{tabular}[b]{l}
 %	  \includegraphics[height=1.3em]{images/IBM-Logo.png}\\
 %	 \end{tabular}
 %	}
 	
 	\newcommand{\emptyHeader}{
 	\fancyhead{}
 	\fancyfoot{}
 	\renewcommand{\headrulewidth}{0pt}
 	}

 }
\fancypagestyle{plain}{
	 	\newcommand{\clearHeader}{\fancyhead{}\fancyfoot{}}
	\fancyhead{}
	\fancyfoot{}  	
 	\fancyfoot[C]{\thepage}
	\fancyhead[L]{\includegraphics[height=0.65cm]{images/dhbw-Logo.png}}
 	
 
 	\fancyhead[C]
 	{
 	 \begin{tabular}[b]{l}
 	  LeMon
 	 \end{tabular}
 	}

 %	\fancyhead[R]
 %	{
 %	 \begin{tabular}[b]{l}
 %	  \includegraphics[height=1.3em]{images/IBM-Logo.png}\\
 %	 \end{tabular}
 %	}
 	
 	\newcommand{\emptyHeader}{
 	\fancyhead{}
 	\fancyfoot{}
 	\renewcommand{\headrulewidth}{0pt}
 	}

}
\fancypagestyle{abstract}{
	 	\newcommand{\clearHeader}{\fancyhead{}\fancyfoot{}}
	\fancyhead{}
	\fancyfoot{}  	
 	\fancyfoot[C]{\thepage}
	\fancyhead[L]{\includegraphics[height=0.65cm]{images/dhbw-Logo.png}}
 	
 
 	\fancyhead[C]
 	{
 	 \begin{tabular}[b]{l}
 	  LeMon
 	 \end{tabular}
 	}

 %	\fancyhead[R]
 %	{
 %	 \begin{tabular}[b]{l}
 %	  \includegraphics[height=1.3em]{images/IBM-Logo.png}\\
 %	 \end{tabular}
 %	}
 	
 	\newcommand{\emptyHeader}{
 	\fancyhead{}
 	\fancyfoot{}
 	\renewcommand{\headrulewidth}{0pt}
 	}

}
%Abkürzungen
%
% Abkürzungen --> referenz, name, beschreibung
% Aufruf mit \gls{...} oder Kurzform mit \acrshort{...}

\newacronym{DHBW}{DHBW}{Duale Hochschule Baden Württemberg}
\newacronym{HDM}{HDM}{Hochschule der Medien}
\newacronym{LeMon}{LeMon}{Lecture Monitoring}
\newacronym{HTML}{HTML}{Hypertext Markup Language}
\newacronym{SQL}{SQL}{Standard Query Language}
\newacronym{Ajax}{Ajax}{Asynchronous JavaScript and XML}

\newglossaryentry{PHP}{name={PHP}, first={PHP}, description={Hypertext Preprocessor, eine Skriptsprache zur serverseitigen Programmierung und Verarbeitung von Webinhalten}}
\newglossaryentry{adoDB}{name={adoDB}, first={adoDB}, description={Ein Framework zum Kapseln der Zugriffe auf Datenbanken}}
\newglossaryentry{jQuery}{name={jQuery}, first={jQuery}, description={Eine JavaScript-Bibliothek zum Manipulieren von browserübergreifenden DOM-Objekten}}
\newglossaryentry{PHPQRCode}{name={PHPQRCode}, first={PHPQRCode}, description={Eine PHP Library zur Erzeugung von 2-dimensionalen Barcodes (sogenannten QR-Codes)}}
\newglossaryentry{pChart}{name={pChart}, first={pChart}, description={Eine PHP Library zur Erzeugung von Diagrammen}}
\newglossaryentry{JavaScript}{name={JavaScript}, first={JavaScript}, description={JavaScript ist eine Script-Sprache zur Clientseitigen Programmierung}}
\newglossaryentry{Bootstrap}{name={Bootstrap}, first={Bootstrap}, description={Bootstrap ist eine JavaScript Bibliothek}}

\newglossaryentry{IDE}{name={IDE}, first={Integrated development environment
(IDE)}, description={Integrated development environment / Integrierte
Entwicklungs-Umgebung}}



% 	\newcommand{\clearHeader}{\fancyhead{}\fancyfoot{}}
	\fancyhead{}
	\fancyfoot{}  	
 	\fancyfoot[C]{\thepage}
	\fancyhead[L]{\includegraphics[height=0.65cm]{images/dhbw-Logo.png}}
 	
 
 	\fancyhead[C]
 	{
 	 \begin{tabular}[b]{l}
 	  LeMon
 	 \end{tabular}
 	}

 %	\fancyhead[R]
 %	{
 %	 \begin{tabular}[b]{l}
 %	  \includegraphics[height=1.3em]{images/IBM-Logo.png}\\
 %	 \end{tabular}
 %	}
 	
 	\newcommand{\emptyHeader}{
 	\fancyhead{}
 	\fancyfoot{}
 	\renewcommand{\headrulewidth}{0pt}
 	}


%Glossar-Befehle anschalten
\makeglossaries

%%%%%%%%%%%%%%%%%%%%%%%%%%%%%%%%%%%%%%%%%%%%%%%%%%%%%%%%%%%%%%%%%%%%%%%%%%%%%%%%


\begin{document} 

\cleardoublepage{}
 
% Deckblatt 
\begin{titlepage}
\begin{longtable}{p{.55\textwidth} p{.85\textwidth}}
%{\includegraphics[height=2.6cm]{images/IBM-Logo.png}} &
{\includegraphics[height=2.6cm]{images/dhbw-Logo.png}}
\end{longtable}
\enlargethispage{20mm}
\begin{center}
\vspace*{8mm} {\LARGE \bf \titel }\\
\vspace*{3mm} {\Large\bf \subtitel} \\
\vspace*{12mm} {\large \bf \arbeit}\\
\vspace*{12mm} {\large \bf \semester}\\
  
\vspace*{12mm} Studiengang  \textit{\studiengang\\}
\vspace*{5mm} Vertiefungsrichtung \textit{\vertiefung\\}
\vspace*{9mm} an der \linebreak Dualen Hochschule Baden-Württemberg / \dhbw\\
\vspace*{12mm} von\\ 
\vspace*{3mm} {\large\bf \autor}\\
\vspace*{12mm} \datumAbgabe\\
\end{center}
\vfill
\begin{spacing}{1.2}
\begin{tabbing}
mmmmmmmmmmmmmmmnmmmmmmmmmmmm \= \kill
\textbf{Bearbeitungszeitraum} \> \zeitraum\\
\textbf{Matrikelnummer Ephraim Petry, Kurs} \> \martrikelnrEphraim, \kurs\\
\textbf{Matrikelnummer Nick Herrmannsdörfer, Kurs}\> \martrikelnrNick, \kurs\\
\textbf{Betreuer der Praxisphase} \> \betreuer\\
\end{tabbing}
\end{spacing}
\end{titlepage}
\renewcommand{\abstractname}{Zusammenfassung}
\begin{abstract} 
 
Im Rahmen dieser Studienarbeit sollte ein System erstellt werden, mit dem der Lernfortschritt und
das Verständnis  der Studenten an der \gls{HDM} für die in der Vorlesung vermittelten Inhalte
abgeprüft werden sollte. Dieses System sollte \gls{LeMon} heißen.

Das System sollte dabei den kompletten Prozess von der Erstellung einzelner Übungsblätter durch den
Dozenten, das Management der Fragen für die Übungsbläter, sowie das Freischalten der Übungsblätter
zur entsprechenden Zeit beinhalten.
Desweiteren sollte der Dozent den Studenten auf einfache Art und Weise das elektronische Übungsblatt
zugänglich machen. Hierzu sollte bevorzugt ein QR-Code eingesetzt werden, den die Studenten mit
ihrem Smartphone und einer entsprechenden App sofort einlesen und zum Übungsblatt gelangen konnten.
\newline Am Ende sollte der Dozent auf einfache Weise die Möglichkeit haben, die Ergebnisse wenige
Sekunden nach Ablauf der Zeit anschaulich angezeigt zu bekommen und diese mit den Studenten
durchsprechen.

Zur Umsetzung dieses Projektes sollte eng mit einer Gruppe mit einem ähnlichen Projekt
zusammengearbeitet werden, die jedoch eine etwas andere Art von Übungsblättern für den Einsatz von
Aufgaben innerhalb oder auch außerhalb der Vorlesung erstellen sollte. Soweit möglich sollten dabei
Synergien genutzt und eine gemeinsame Oberfläche für beide Projekte erstellt werden.




 
 \end{abstract} 
\glsresetall
\renewcommand{\baselinestretch}{1.50}\normalsize
%\chapter*{Vorwort}
Dies ist die Studienarbeit von Ephraim Petry und Nick Herrmansdörfer. Sie ist bestandteil des
dritten Studienjahres des Bachelor-Studiums an der \gls{DHBW} Stuttgart.

Diese Studienarbeit sollte nicht nur den Zweck der Absolvierung einer Studienarbeit erfüllen. Ein
besonderes Augenmerk lag auch darauf, dass bestehende oberflächliche Wissen in bestimmten
Technologien in als praktisches Projekt umzusetzen und zu vertiefen.

\tableofcontents 
\glsresetall 
\chapter{Einleitung}


Hier könnte eine Einleitung stehen \emph{Hervorgehoben} 
 
\section{Überschrift abc}
 Abcdef ghijklmno pqrstuvw xyz äöü\ldots \gls{DHBW}
%  
% \section{Überschrift abcd}\label{sec:Ueberschriftabcd}
% Abcdef ghijklmno pqrstuvw xyz ÄÖÜ\ldots 
% 
% \prettyref{sec:Ueberschriftabcd}
% 
% 
% \begin{figure} [!htb]
% 	\begin{center}
% 		  \begin{tabular}{@{}r@{}}
% 		{\includegraphics[width=14cm]{images/dhbw-Logo.png}}\\
% 		\footnotesize\sffamily\textbf{Quelle:} DHBW-Homepage
% 	  		   \cite{LiteraturEintrag1}
%  	 	 \end{tabular} 	
% 		\caption{Das DHBW-Logo}
% 		\label{fig:DHBWLogo}
% 	\end{center} 
% \end{figure}
\chapter{Use Cases}
As MapReduce is an important paradigm for solving big and costly problems today, there are really
many use cases. Almost every problem that can be divided into small pieces and processed independent
can be distributed and calculated in parallel by following the MapReduce paradigm. Not only
scientific calculation, but also today’s companies are trying to solve their problems by using Map
Reduce.

The following shows an exemplary list of problems that can be efficiently solved by MapReduce:
\begin{singlespacing}
	\begin{itemize}
	  \item Factorization of (big) Integers (Quelle: https://code.google.com/p/mapreduce-integer-factorization/) 
	  \item Matrix factorization (Quelle: http://www.redaktion.tu-berlin.de/fileadmin/fg131/Publikation/Papers/sys024-schelter.pdf) 
	  \item Multiplying terabit integers (Quelle: https://people.apache.org/~szetszwo/ssmr20110430.pdf)
	  \item Fourier transformations
	  \item Genetic analysis (Quelle: http://genome.cshlp.org/content/20/9/1297.long) 
	  \item Search engine indexing
	  \item Big sort problems (Quelle: http://courses.cs.washington.edu/courses/cse490h/08au/lectures/algorithms.pdf) 
	  \item Users' behavior analytics
	  \item Group text documents into topically related groups
	  \item ETL and Data Mining
	  \item Friend finder
	  \item Defending E-Mail-Spam (Yahoo Mail: http://readwrite.com/2010/05/24/map-reduce-yahoo-mail#awesm=~ozyRktADsuJsXT) 
	  \item Shortest Path (Quelle: http://courses.cs.washington.edu/courses/cse490h/08au/lectures/algorithms.pdf)  
	  \item Word count (Quelle: http://courses.cs.washington.edu/courses/cse490h/08au/lectures/algorithms.pdf) 
	  \item Weather prediction
	\end{itemize}
\end{singlespacing}

Some of these use cases are more scientific, others are for commercial use. Especially while data
warehouses are rising, there is a need for calculating big data. When data is filled into the data
ware house, there are often calculations to be done like calculating aggregations. This can be
efficiently done by using Map Reduce. When data is already filled inside the data ware house, there
are several techniques for data mining. As many parts of data mining can be split into pieces and
calculated separately, it is very efficient to take Map Reduce algorithms to do so.

As MapReduce processing is computed by large clusters, nowadays with many cloud providers it is more
a matter of being interested in the results than having the money to buy an own cluster. Almost all
providers of cloud computing are having resources or hardware optimized solutions for doing
MapReduce. E.g. Amazons Elastic Map Reduce (EMR) is between $0.015 and $0.50 per hour (Quelle:
http://aws.amazon.com/de/elasticmapreduce/pricing/), which isn’t really expensive - especially if it
is for commercial use.

Genetic analysis is an upcoming trend. Today the costs for a complete genetic analysis are below
\$1000. By doing this analysis on public cloud infrastructure and doing the analysis with the Map
Reduce paradigm, these costs could be significantly reduced. Maybe in some years, it will be a
normal thing to analyses the genomes of a person? - The technic is ready to do so. The bigger
question is: Are the people willing to do that?

\begin{figure} [!htb]
	\begin{center}
		  \begin{tabular}{@{}r@{}}
		{\includegraphics[width=14cm]{images/genome.jpg}}\\
		\footnotesize\sffamily\textbf{Source:} \url{http://www.medcitynews.com}
	  		   \cite{genomePicture} 
 	 	 \end{tabular}
		\caption{Picture of a genome}
		\label{fig:genome}
	\end{center}
\end{figure}\FloatBarrier

Factorization of big integers is a very interesting but also a dangerous use case. Many security and
cryptographic methods are based on the problem that it is not possible to efficiently factorize big
integers. Map Reduce is a way to solve some factorization problems. But if someone could find a
really efficient way (much more efficient than Map Reduce), this would lead to new challenges in
cryptography.

Search engine companies like Google do have a big need for indexing huge amount of data to give
their users the best and most relevant search results in the shortest possible time. It is hardly
surprising that Google as the market leader and inventor of MapReduce uses MapReduce to do that job
while it is highly parallelizable. Google also uses the same techniques for indexing Emails in
Gmail.

Friend finders were a revolutionary new thing for social networks. Companies.  With Facebook’s
friend finder leading the way big social networks implemented better search engines for people who
are searching for their friends. A large amount of data has to be analyzed to give good suggestions
for possible friends. Indexing profiles and searching for people can be highly parallelized.
Facebook implemented its friend finder by using the MapReduce pattern.

\begin{figure} [!htb]
	\begin{center}
		  \begin{tabular}{@{}r@{}}
		{\includegraphics[width=14cm]{images/FacebookfriendFinder.png}}\\
 	 	 \end{tabular}
		\caption{Facebooks friend finder}
		\label{fig:friendFinder}
	\end{center}
\end{figure}

\chapter{Fazit und Ausblick}

 

% Anhang
\chapter{Anhang} 
\clearpage
\pagenumbering{Roman}

%Glossar ausgeben
\printglossary[ 
	%%style=indexgroup,
	title=Glossar,
	toctitle=Glossar]
    
%   separates Abkürzungsverzeichnis ausgeben  
%   \printglossary[
%   	style=indexgroup,
%   	type=\acronymtype, 
%   	title =Abkürzungsverzeichnis,
%		toctitle=Abkürzungsverzeichnis]

% Abbildungsverzeichnis
\clearpage

%Quellcode-Anhang
%\cleardoublepage
%\phantomsection \label{listofCode}
%\addcontentsline{toc}{section}{Quellcode1}
%\clearpage


\phantomsection \label{listoffig}
\addcontentsline{toc}{section}{Abbildungsverzeichnis}
\listoffigures

%Tabellenverzeichnis
\cleardoublepage
\phantomsection \label{listoftab}
\addcontentsline{toc}{section}{Tabellenverzeichnis}
\listoftables 
 
% Quellcodeverzeichnis
\cleardoublepage
\phantomsection \label{listoflist} 
\addcontentsline{toc}{section}{Quellcodeverzeichnis}
\lstlistoflistings

% im Literaturverzeichnis alles anzeigen, auch die die man nicht verwendet hat 
% \nocite{*} 
% Literaturverzeichnis (durch BibTeX erstellt)

\interlinepenalty=10000 %Seitenumbrüche in einzelnen Einträgen im
% Literaturverzeichnis vermeiden
\bibliography{literaturverzeichnis}{}
\clearpage

% leere Seite am Ende einfügen
%\thispagestyle{empty}
%\textbf{ }
%\cleardoublepage


\end{document} 